%==============================================================================
\section{Conclusions and further work}\label{sec:concl}
%==============================================================================

\subsection{What this report establishes}\label{ssec:concl-summary}

The equivalence between the SSP and a walk through magazine configurations organises
everything above. It yields a structural theory of the standard heuristic and a family
of exact methods, and one quantity --- the coverage bound --- governs both.

\paragraph{On the structural side.}
The two-phase heuristic admits two readings, and separating them is a prerequisite for
saying anything correct about it (Section~\ref{ssec:frames}). Within that setting:

\begin{itemize}[topsep=3pt,itemsep=2pt]
  \item The optimum equals the cheapest walk over \emph{all} feasible groupings
  (Theorem~\ref{thm:grouping}), so the best value the two-phase construction can attain
  differs from the optimum exactly because of the minimum-cardinality restriction
  (Corollary~\ref{cor:gap}). This bounds the envelope $H$ of
  Definition~\ref{def:frames}, not the output of a particular tie-breaking
  implementation.
  \item Zero-gap classes exist for every capacity: two groups, or a tool set barely
  exceeding the magazine, or a small optimum (Corollaries~\ref{cor:zerogap} and
  \ref{cor:smallZ}).
  \item The ratio is bounded for every instance by
  $\min(b,\Kst-1,q)$ (Theorem~\ref{thm:uncond}), which tightens automatically on
  instances with few groups or a small tool excess.
  \item At three groups the walk cost has a closed form
  (Proposition~\ref{prop:hk3}), and the gap obeys a counting bound
  (Proposition~\ref{prop:k3}) that is attained by explicit witnesses. The natural
  guess $\mathrm{gap}\le\Kst-2$ is false for the walk value
  (Proposition~\ref{prop:refute}).
  \item For the method as actually specified, positive gaps are rare in the censuses,
  equal to one on every connected instance found, and never occur at $b=2$
  (Observations~\ref{obs:rings}--\ref{obs:census}). Unboundedness survives, but with a
  different witness family from the one the walk model suggests. That family is
  \emph{disconnected}, and its unboundedness is a conjecture supported at $g=1,2$ rather
  than a theorem (Observation~\ref{obs:unbounded-true} and
  Conjecture~\ref{conj:unbounded}). Whether the gap is bounded on \emph{connected}
  instances is open.
  \item The circuit clutter of the grouping system computes $\Kst$
  (Proposition~\ref{prop:chromK}) but does not determine the gap
  (Proposition~\ref{prop:noclutter}), while blocking duality gives a first
  non-integral family for the grouping relaxation
  (Proposition~\ref{prop:oddring}).
  \item A per-changeover setup cost interpolates between the SSP and the JGP, with the
  transition confined to an interval whose width is the gap itself
  (Theorem~\ref{thm:collapse} and Lemma~\ref{lem:lowrho}). This answers a question
  posed by \citet{Burger2015}.
\end{itemize}

\paragraph{On the algorithmic side.}
A branch-and-Benders-cut solver was built whose subproblem is exact and polynomial,
together with position-indexed formulations carrying a fixed polynomial row set, among
which the position--transition model is a configuration relaxation that can provably
improve on the coverage bound (Proposition~\ref{prop:ptf}), although the campaign finds
it doing so on three runs out of $233$.

\paragraph{On the empirical side.}
A campaign over $1{,}421$ instances and twelve configurations, under a single one-hour
limit, is reported in Section~\ref{sec:comp} and summarised in
Section~\ref{ssec:results-summary}. Four of its measurements --- the incumbent at the
time limit, the ablation, the split by whether the coverage bound is tight, and the
root relaxation reported by the branch-and-price prototypes --- are independent of one
another and agree. What they agree on is that the obstacle is the bound rather than the
search: nothing that improves the search closes the loose instances, and nothing built
here improves the bound. Two further measurements exist to make that reading
trustworthy rather than to support it, and both do: no two methods that proved the same
optimum disagree on its value, and the reimplemented compact baselines are reported
without a ranking, because the runs that failed are the ones that would have decided it.

Two things this campaign is not should be said plainly here rather than left to
Section~\ref{ssec:limits}. The compact baselines close more instances than
anything built for this report, and they do so while being represented by F4 rather
than the F5 that \citet{Catanzaro2015} recommend --- so the margin is if anything
understated, and no claim of a faster solver is being made. And the branch-and-price
prototypes were run to read their root relaxation, not to compete: their times measure
a Python pricer.

The methods built here are therefore bound-limited rather than
implementation-limited: a negative result about a family of methods rather than about
one solver. Section~\ref{ssec:disc-locality} sets out where a stronger bound could
come from.

\subsection{Further work}\label{ssec:further}

\paragraph{The exact worst case of the heuristic.}
Section~\ref{ssec:law} settles the extremal behaviour of the walk value at three
groups. For the value $H$ of Definition~\ref{def:frames}, the question is open and is now the more
interesting one, because Section~\ref{ssec:ktnsgap} shows that the two differ
substantially. Two concrete questions: is the heuristic gap ever larger than one on a
connected instance, and does $\mathrm{gap}\le\Kst-2$ hold for it? Directed searches
across $b\in\rset{3,4,5}$ found no counterexample to either.
Problem~\ref{prob:ktns-closes} is the structural question underneath both.

\paragraph{The polyhedral questions.}
One of the two questions from Section~\ref{ssec:clutter} is now answered.
Observation~\ref{obs:starobstruction} enumerates the whole $b=3$ family with two tools
per job on at most six tools: half of it has a non-integral covering relaxation, eleven
of those instances are bipartite, and the smallest obstruction is a three-job star
rather than a ring. What the star and the odd ring share is odd structure in the
compatibility hypergraph on \emph{jobs}, not circular structure among the tools, and
characterising exactly which such structures obstruct integrality is the question that
replaces the one that has been closed. The second question is unchanged: which facets of
the position--transition polytope explain the three runs where its bound exceeds $q$.

\paragraph{A bound that reaches into the loose half.}
This is the direction the campaign points to most strongly, and
Section~\ref{ssec:disc-locality} takes the first step along it: the arc-supported
information the solvers carry today is exhausted, and the window family of
inequality~\eqref{eq:window} has roughly three and a half times as much room at its
\emph{ceiling over integral schedules} (Table~\ref{tab:ceilings}). A ceiling is not a
bound. That family has now been measured inside the master
and it delivers nothing: the root relaxation is unchanged on all $37$ instances tested
(Observation~\ref{obs:windowroot}). The robust encoding is the reason. Counting the
tools an integral schedule would need inside a window requires tying the row to the
column variables, which puts the cut's dual into the pricing problem and forfeits the one
property that made it cheap. Whether a non-robust encoding pays for itself is open, and
so is whether the rows help inside the tree despite leaving the root alone. Two further
families suggest themselves from the same
reading: chaining PTF per configuration rather than per tool, which is what confines
its advantage to three runs in Section~\ref{ssec:results-bnp}, and disaggregating the
Benders master's single $\theta$ into one variable per position, which is the standard
strengthening the ablation of Section~\ref{ssec:results-bbc} never tested.

\paragraph{Making PTF's advantage large.}
Section~\ref{ssec:results-bnp} finds the position--transition relaxation above the
coverage bound on three runs of $233$ --- by one unit on two of them and by $0.0134$ on
the third. The construction is therefore
sound and, as it stands, of no practical use. Problem~\ref{prob:ptf-gap} asks whether
the excess can be made to grow; the polyhedral question underneath it is which facets
of the position--transition polytope are responsible for the three cases where it does
exceed $q$. Independently, the prototypes are written in Python with a Python pricer,
and a compiled pricer would allow the formulations to be measured on instances larger
than nine jobs, which is where they currently stop.

\paragraph{Non-uniform tool costs.}
Section~\ref{ssec:bbc} notes that the Benders subproblem is the loading linear program
rather than the KTNS rule, and that \citet{Crama1994} show the linear program remains
exact when tools have different setup times while the rule does not. The solver
therefore extends to that regime unchanged. Testing it there would broaden the scope of
the method beyond the uniform problem that the whole SSP literature assumes.

\paragraph{Beyond the SSP.}
Nothing in the position--transition construction is specific to tool switching except
the transition metric. Position-transition columns with per-element consistency rows
apply to any generalised travelling-salesman problem with overlapping clusters
\citep{Pop2024}. Whether the relaxation strength observed here persists in that
generality is open.

\paragraph{Learning-augmented methods.}
Machine learning for the exact side is an active area, covering cut selection
\citep{Tang2020,DezaKhalil2023} and column and branching selection
\citep{Bengio2021,Gasse2019,Morabit2021}. Its established role is to \emph{choose}
among cuts or columns one already knows how to generate, not to make them stronger.
Section~\ref{sec:disc} bounds how much that can help here: what defeats the solver is
a bound gap, and a selector cannot select a strong cut that does not exist. If a
learned component is to help this problem, the more promising target is column
selection for the position-indexed pricers \citep{Morabit2021}, whose coupled-pair
oracle is the measured bottleneck, rather than cut selection for the Benders master.
On the pricing side the corresponding non-learned upgrade is the automatic dual
stabilisation of \citet{Pessoa2018}.
